Follow this procedure to quickly reduce \instname\ demo data. We
assume that the \edps\ Graphic User Interface and the MUSE pipeline
and the demo data are installed in your system. For general
instructions on how to install \edps\ \edpsgui\ the pipeline and all
the needed components, please visit
\href{https://www.eso.org/sci/software/pipe_aem_main.html}{this page}.


\subsection{Setting the workflow}

\begin{enumerate}[a)]
\item blabla
\end{enumerate}

\begin{enumerate}
\item After installation and while still in the installation virtual environment (if not, activate it \verb| [again](../quick/installation.md#installation-steps)|), start the \edpsgui\ dashboard by typing:

\begin{verbatim}
  edps-gui
\end{verbatim}

The \edpsgui\ will start in a browser window. 

\begin{verbatim}
   :alt: dashboard
   :name: fig_dashboard
\end{verbatim}

\item Optionally, before starting \edps\ one can specify new settings by pressing Help --> Settings (see \verb|[here](../edpsgui/settings_.md)|)

\begin{verbatim}
   :alt: help_settings
   :name: fig_help_settings
\end{verbatim}

\item On the browser window with the \edpsgui\ dashboard, press the button \texttt{Start EDPS}.

\begin{verbatim}
   :alt: start_edps
   :name: fig_start_edps
\end{verbatim}

\begin{verbatim}
```{figure} figures/select_workflow.jpg
:alt: wkf
:name: fig-wkf

The EDPS Dashboard (Graphic User Interface) with the MUSE workflow loaded.	
\end{verbatim}

\item Choose the \wkfname\ pipeline workflow from the list in the `workflows` field. The workflows offered in this selector depend on
the installed pipelines.
The graphic workflow representation will appear as in
\verb|{numref}`fig-wkf`|.

\end{enumerate}

\subsection{Selecting the input data}

\begin{enumerate}[a)]
\item  Press the button \texttt{Raw Data}.
\begin{verbatim}
   :alt: raw_data
   :name: fig_raw_data
\end{verbatim}

\item  Press \texttt{Select Inputs}. A selection window will appear that allows to select data that are stored on a local disk.

\begin{verbatim}
   :alt: select_input
   :name: fig_select_input
\end{verbatim}

\item (Optional). Select the reduction target, configure the workflow parameter and specify the association preferences. 
Thise steps are optional. For more information see \verb|[here](configure_dataset_.md)|.

\item Press \texttt{Create Datasets}. A list of datasets appears, one line for each set of science data.

\begin{verbatim}
   :alt: create_datasets
   :name: fig_create_datasets
\end{verbatim}


\item Choose the datasets that should be processed 
\begin{verbatim}
   :alt: send
   :name: fig_send
\end{verbatim}
and send them to the data reduction queue by pressing \texttt{Submit to Reduction Queue}. Note that this action does not start the 
reduction automatically.

\end{enumerate}

\subsection{Start the reduction}
To start the reduction of the datasets that are submitted to the reduction queue:
\begin{enumerate}[a)]
\item  Press \texttt{Reduction Queue}.

\begin{verbatim}
   :alt: processing_queue
   :name: fig_processing_queue
\end{verbatim}

\item  (Optional). Configure the workflow and recipe parameters by pressing the wheel 
button \verb|[](figures/configure\_dataset.jpg)| to open the configuration editor.

\item  Press the \texttt{Reduce} button. The selected data will now be processed with the default parameters (unless edited in the previous step).

\begin{verbatim}
   :alt: reduce
   :name: fig_reduce
\end{verbatim}

\end{enumerate}


Congratulations! You reduced your first data with \edpsgui\. All the
reduced data are saved in the EDPS\_data directory specified when
executing \edpsgui\ for the first time.
